\documentclass[10pt,a4paper]{amsart}
\usepackage[margin=19mm]{geometry}
\usepackage{amsmath,amssymb,mathtools}
\usepackage[scr=rsfs,cal=euler]{mathalfa}
\usepackage{xfrac}
\usepackage[unicode,hidelinks,bookmarks=false]{hyperref}
\newcommand{\ie}{i.e.\ }
\newcommand{\cf}{cf.\ }
\newcommand{\resp}{respectively\ }
\newcommand{\Subsection}{\subsection}
\providecommand{\nobreakdash}{\nobreak\hbox{-}}
\setlength{\emergencystretch}{2em}
\multlinegap0pt
\usepackage{bm}
\usepackage{bbm}
\usepackage{dsfont}
\usepackage{xspace}
\usepackage{cancel}
\usepackage{mathtools}
\usepackage{amsthm}
\makeatletter
\@ifundefined{Theorem}{\newtheorem{Theorem}{Theorem}}{}
\@ifundefined{Lemma}{\newtheorem{Lemma}[Theorem]{Lemma}}{}
\makeatother
\title[Control and stabilization of cascade coupled systems: applic]{Control and stabilization of cascade coupled systems: application to a 1-d heat and wave coupled system}
\author{Lucas Davron, Pierre Lissy, Swann Marx}
\date{Source: arXiv:2603.09829v1 (CC BY 4.0)}
\begin{document}
\maketitle

\noindent\emph{Source and licence.} Control and stabilization of cascade coupled systems: application to a 1-d heat and wave coupled system. Version arXiv:2603.09829v1 (CC BY 4.0), distributed under the Creative Commons Attribution 4.0 International licence, as recorded in the arXiv licence metadata for this exact version and shown on the abstract page https://arxiv.org/abs/2603.09829v1. Full rights notice: licenses/arxiv-2603.09829-CC-BY-4.0.txt.\par\smallskip

\noindent\textit{Editorial scope.} Counted unit: the Lemma whose source label is \texttt{lem:spectral\_asymp} in the article (source0.tex lines 1218--1232, source label \texttt{lem:spectral\_asymp}), delivered together with the article's own complete original proof of it (source0.tex lines 1233--1296, one \texttt{proof} environment of 64 lines). No mathematics is written, repaired or re-derived: statement and proof are the article's own text line for line. Required context retained uncounted: none. Every definition, display and label of those lines is retained unchanged. Omitted from this package and not redistributed: 1--1217, 1297--1690. Nothing inside a retained excerpt is omitted or altered: every retained body is delivered exactly as the article's lines are written, so no omission receipt of any kind accompanies this package. Citations: the retained excerpts carry no citation command at all, so no bibliography entry of the article is transcribed and no reference is redistributed. Reference adaptation: none was needed and none was made; every reference inside the delivered unit resolves inside it, so no unresolved reference or citation is produced. Printed slips kept literal: no printed word, symbol, hypothesis or number of the counted statement or of its proof was altered. Layout and class: the article's own class is \texttt{article} with packages including amsmath, manfnt, amssymb, graphicx, geometry, amsthm, xcolor, inputenc, pst-node, tikz-cd, hyperref, bbold, pdfpages, biblatex; the wrapper is the lane's pinned \texttt{amsart} editorial wrapper at 10pt on A4, which restates the theorem-like environments the retained text needs and adds the article's own macro definitions that the retained text needs (none), with extra wrapper packages (bm, bbm, dsfont, xspace, cancel, mathtools). The delivered numbering is the unit's own. Assets: no retained span contains a figure environment, a graphics-inclusion command, a PDF-page inclusion, a table or a TikZ picture; no image file is copied into this package, the package declares no asset dependency and no image file is redistributed with it. Licence: version arXiv:2603.09829v1 of this article is distributed under the Creative Commons Attribution 4.0 International licence, as recorded in the arXiv licence metadata for this exact version and shown on the abstract page https://arxiv.org/abs/2603.09829v1. A full rights notice is delivered at licenses/arxiv-2603.09829-CC-BY-4.0.txt. Characters: every retained excerpt is pure ASCII, so the delivered engine, XeLaTeX with its default Latin Modern text font, reports no missing character.
\begin{Lemma}\label{lem:spectral_asymp}
The eigenvalues of $-A_\alpha$ satisfy 
\begin{equation}\label{eq:spectral_genuine}
\forall \alpha > 0,\quad \forall j \geq 1,\quad (j-1)^2\pi^2 < \lambda_{j,\alpha} < \left( j - \frac{1}{2}\right)^2 \pi^2.
\end{equation}
Moreover,
\[
 \sqrt{\lambda_{1,\alpha}} = \sqrt{\lambda_{1,0}} + \alpha^{1/2} - \frac{1}{6} \alpha^{3/2} + \frac{2}{45} \alpha^{5/2} +  O(\alpha^3),\quad \alpha \to 0^+.
 \]
and 
\[
\sqrt{\lambda_{j,\alpha}} = \sqrt{\lambda_{j,0}} + \frac{\alpha}{\sqrt{\lambda_{j,0}}} - \frac{\alpha^2}{\lambda_{j,0}^{3/2}} + \frac{o_j(\alpha^2)}{\lambda_{j,0}^{3/2}}, \quad \alpha \to 0^+,\quad j \geq 2,
\]
where the term $o_j(\alpha^2)$ depends on $j$ and $\alpha > 0$ but has an implicit constant that is uniform in $j \geq 2$ with respect to $\alpha \to 0^+$\footnote{To simplify, we will say that $o_j$ is uniform in $j$.}. 
\end{Lemma}

\begin{proof}
An elementary computation shows that if $\lambda > 0$ is an eigenvalue of $-A_\alpha$, then a corresponding eigenvector is $\cos(\sqrt{\lambda}x)$ and we have the equation 
\begin{equation}\label{eq:vp_robin}
    - \sqrt{\lambda} \sin(\sqrt{\lambda}) + \alpha \cos(\sqrt{\lambda}) = 0. 
\end{equation}
One readily verifies that $\sqrt{\lambda}$ cannot be an integer multiple of $\pi$, hence $(j-1) \pi < \sqrt{\lambda} < j\pi$ for some $j \geq 1$. By sign considerations, one further sees that $(j-1) \pi < \sqrt{\lambda} < (j-1/2)\pi$. Conversely, for all $j \geq 1$, the intermediate value theorem yields a unique solution to \eqref{eq:vp_robin} in $(j-1) \pi < \sqrt{\lambda} < (j-1/2)\pi$. Therefore, the eigenvalues of $-A_\alpha$ are simple and $(j-1) \pi < \sqrt{\lambda_{j,\alpha}} < (j-1/2)\pi$. 

We first show the asymptotics on $\sqrt{\lambda_{1,\alpha}}$. To match later notations, set $x_{1,\alpha}=\sqrt{\lambda_{1,\alpha}}$. We have $\alpha = f(x_{1,\alpha})$ with $f(x) = x \tan x$, which is a smooth bijection from $(0,\pi/2)$ to $(0,\infty)$. Thus $x_{1,\alpha} \to 0^+$ as $\alpha \to 0^+$. We have 
\[
x_{1,\alpha} = a\alpha^{1/2} +b\alpha^{3/2} +c \alpha^{5/2} + O(\alpha^3) \Longleftrightarrow \alpha = f\left( a\alpha^{1/2} +b\alpha^{3/2} +c \alpha^{5/2} + O(\alpha^3) \right),
\]
and since   
\[
f(x) = x^2 + \frac{1}{3}x^4 + \frac{2}{15} x^6 + O(x^8) ,\quad x \to 0^+,
\]
we deduce 
\[
f\left( a\alpha^{1/2} +b\alpha^{3/2} +c \alpha^{5/2} + O(\alpha^3) \right) = a^2 \alpha + \left(2ab + \frac{a^4}{3}\right)\alpha^2 + \left(2ac + \frac{4a^3b}{3} + \frac{2a^6}{15}\right) \alpha^3 + O(\alpha^{7/2}).
\]
The system 
\[
\left\lbrace \begin{array}{lcc}
a^2 &=& 1,\\
\displaystyle 2ab + \frac{a^4}{3} &=& 0, \\
\displaystyle 2ac + \frac{4a^3b}{3} + \frac{2a^6}{15} &=& 0, \\
\end{array}\right.
\]
has the solution $(a,b,c) = (1,-1/6,2/45)$, hence the asymptotics on $x_{1,\alpha}$. 


Now fix $j \geq 2$, the equation \eqref{eq:vp_robin} rewrites in the new variable $x_{j,\alpha} := \sqrt{\lambda_{j,\alpha}} - (j-1)\pi \in (0,\pi/2)$ as
\[
\left( \frac{x_{j,\alpha}}{\sqrt{\lambda_{j,0}}} + 1 \right) \tan(x_{j,\alpha}) = \frac{\alpha}{\sqrt{\lambda_{j,0}}},
\]
owing to the $\pi$-periodicity of $\tan$. Put 
\[
f_j(x) := \left( \frac{x}{\sqrt{\lambda_{j,0}}} + 1 \right) \tan(x),
\]
which is smooth on $(0,\pi/2)$ and increases from $f_j(0^+) = 0$ to $f_j(\pi/2^-) = +\infty$, hence for fixed $j \geq 2$ we have $x_{j,\alpha} \to 0^+$ as $\alpha \to 0^+$. We adapt the above method, for fixed constants $a,b,c,d,$ we have
\begin{equation}\label{xjep}
x_{j,\alpha} = a \frac{\alpha}{\sqrt{\lambda_{j,0}}} + b \frac{\alpha^2}{\sqrt{\lambda_{j,0}}} + c \frac{\alpha^2}{\lambda_{j,0}} + d \frac{\alpha^2}{\lambda_{j,0}^{3/2}} + \frac{O_j(\alpha^3)}{\lambda_{j,0}^{3/2}},
\end{equation}
as $\alpha \to 0^+$ and uniformly in $j \geq 2$, if and only if 
\[
\frac{\alpha}{\sqrt{\lambda_{j,0}}} = f_j\left(  a \frac{\alpha}{\sqrt{\lambda_{j,0}}} + b \frac{\alpha^2}{\sqrt{\lambda_{j,0}}} + c \frac{\alpha^2}{\lambda_{j,0}} + d \frac{\alpha^2}{\lambda_{j,0}^{3/2}} + \frac{O_j(\alpha^3)}{\lambda_{j,0}^{3/2}} \right),
\]
as $\alpha \to 0^+$ and uniformly in $j \geq 2$. From the definition of $f_j$, we see that
\[
f_j(x) = \left( \frac{x}{\sqrt{\lambda_{j,0}}} + 1 \right)\left(
x + O(x^3)\right),
\qquad x \to 0^+,
\]
where the $O$ is independent on $j$,
so that
\[
f_j(x_{j,\alpha})
=
x_{j,\alpha}+\frac{x_{j,\alpha}^2}{\sqrt{\lambda_{j,0}}}
+O_j(x_{j,\alpha}^3),
\qquad  \alpha \to 0^+,
\]
where we used that $x_{j,\alpha}/\sqrt{\lambda_{j,0}}$ is bounded from above independently on $j$,  so that  $O_j$ is uniform in $j$ or $\alpha$.
Since $f_j(x_{j,\alpha})={\alpha}/{\sqrt{\lambda_{j,0}}}$, and using \eqref{xjep}, we see that the claimed expansion for $x_{j,\alpha}$ is valid with $(a,b,c,d) = (1,0,0,-1)$, for some $o_j$ that is uniform in $j$, as it can be easily verified, since developing properly the reminders in the expansions give terms of the form $O_j(\alpha^n)/(\sqrt{\lambda_{j,0}})^m$, with $O_j$ uniform in $j$, $n\geqslant 3$ and $m\geqslant 3$.
\end{proof}

\end{document}
